% GENERATED FILE. Edit canonical lessons, then run npm run generate:books.
% canonical-source-hash: fnv1a64:0b9949882d264383
% canonical-lessons: ZH-C454-fuwutai, ZH-C454-zixun, ZH-C454-tushu, ZH-W454-tu-chart, ZH-C454-zhilu, ZH-C454-xuanchuance

\chapter{The Information Desk}
\label{ch:zh-the-information-desk}
\hlchaptermodality{454}

\begin{chapteropening}
\textbf{By the end of this chapter:} \emph{I can find the information desk, make an inquiry, talk about books, give directions and pick up a brochure.}

Complete the last lesson of chapter 454: I can find the information desk, make an inquiry, talk about books, give directions and pick up a brochure.
\end{chapteropening}

\section[fúwùtái]{fúwùtái — service desk; information desk}

\label{lesson:ZH-C454-fuwutai}



\begin{quote}
Before the new one: say the Mandarin for to carry in the hand, then the Mandarin for a luggage tag.
\end{quote}

\subsection*{You'll want to know: fúwùtái}
\emph{fúwùtái} — \textquotedblleft{}service desk; information desk\textquotedblright{}.

\begin{quote}
\emph{qù fúwùtái zěnme zǒu} — How do I get to the information desk?
\end{quote}

Its characters come later; for now it is a word for the ear.

\begin{cousinweb}
\emph{fúwùtái} is built from \emph{fúwù} 'to serve' and \emph{tái} 'a platform', all inherited from Middle Chinese, \emph{bjuwk}, \emph{mjuH} and \emph{doj} (Guangyun, in Baxter's transcription), which glosses the first 'to serve', the second 'affairs' and the last as a high square terrace.
\end{cousinweb}

\begin{grammarlens}[title={What you've built}]
One more word for getting around.
\end{grammarlens}

\subsection*{Your turn}
\begin{itemize}
\raggedright
  \item \emph{Say it:} \emph{fúwùtái}
  \item \emph{Say it:} \emph{fúwùtái}, once more
  \item \emph{Say it:} \emph{qù fúwùtái zěnme zǒu} — How do I get to the information desk?
  \item \emph{Recall:} say the Mandarin for to carry in the hand, then the Mandarin for a luggage tag, then say \emph{fúwùtái} again
\end{itemize}

\begin{tcolorbox}[breakable,colback=teal!4,colframe=teal!35!black,title={Before you move on}]
What does fúwùtái mean? (\textquotedblleft{}Service desk; information desk\textquotedblright{}.) Say it once more.
\end{tcolorbox}
\section[zīxún]{zīxún — to make an inquiry; to ask for advice}

\label{lesson:ZH-C454-zixun}



\begin{quote}
Before the new one: say the Mandarin for a luggage tag, then the Mandarin for an information desk.
\end{quote}

\subsection*{You'll want to know: zīxún}
\emph{zīxún} — \textquotedblleft{}to make an inquiry; to ask for advice\textquotedblright{}.

\begin{quote}
\emph{nǐ kěyǐ qù fúwùtái zīxún} — You can ask at the information desk.
\end{quote}

Its characters come later; for now it is a word for the ear.

\begin{cousinweb}
\emph{zīxún} is built from \emph{zī} 'to consult' and \emph{xún} 'to ask', both inherited from Middle Chinese, \emph{tsij} and \emph{swin} (Guangyun, in Baxter's transcription), which glosses the first 'to take counsel' and the second 'to consult'.
\end{cousinweb}

\begin{grammarlens}[title={What you've built}]
One more word for getting around.
\end{grammarlens}

\subsection*{Your turn}
\begin{itemize}
\raggedright
  \item \emph{Say it:} \emph{zīxún}
  \item \emph{Say it:} \emph{zīxún}, once more
  \item \emph{Say it:} \emph{nǐ kěyǐ qù fúwùtái zīxún} — You can ask at the information desk.
  \item \emph{Recall:} say the Mandarin for a luggage tag, then the Mandarin for an information desk, then say \emph{zīxún} again
\end{itemize}

\begin{tcolorbox}[breakable,colback=teal!4,colframe=teal!35!black,title={Before you move on}]
What does zīxún mean? (\textquotedblleft{}To make an inquiry; to ask for advice\textquotedblright{}.) Say it once more.
\end{tcolorbox}
\section[\texorpdfstring{\zh{图书} (túshū)}{图书 (túshū)}]{\zh{图书} (túshū) — books on the shelves (of a library or bookshop)}

\label{lesson:ZH-C454-tushu}



\begin{quote}
Before the new one: say the Mandarin for an information desk, then the Mandarin for to make an inquiry.
\end{quote}

\subsection*{You'll want to know: \zh{图书}}
\textbf{\zh{图书}} — \emph{túshū} — \textquotedblleft{}books on the shelves (of a library or bookshop)\textquotedblright{}.

\begin{quote}
\textbf{\zh{这家书店的中文图书很多}} \emph{zhè jiā shūdiàn de Zhōngwén túshū hěn duō} — This bookshop has lots of Chinese books.
\end{quote}

The new character, in pieces: \zh{图} — eight strokes: the frame of \zh{国}, and inside it a short stroke falling to the left, a bar that turns down to the left, a long stroke falling to the right, and two dots. The next lesson puts it in your hand.

\begin{cousinweb}
\textbf{\zh{图书}} (\emph{túshū}) is built from \zh{图} 'a picture' and \zh{书} 'a book', both inherited from Middle Chinese, \emph{du} and \emph{syo} (Guangyun, in Baxter's transcription), which glosses the first 'to plan'; the sense 'a picture, a map' is not given there.
\end{cousinweb}

\begin{grammarlens}[title={What you've built}]
One more word for getting around.
\end{grammarlens}

\subsection*{Your turn}
\begin{itemize}
\raggedright
  \item \emph{Say it:} \emph{túshū}
  \item \emph{Say it:} \emph{túshū}, once more
  \item \emph{Say it:} \emph{zhè jiā shūdiàn de Zhōngwén túshū hěn duō} — This bookshop has lots of Chinese books.
  \item \emph{Recall:} say the Mandarin for an information desk, then the Mandarin for to make an inquiry, then say \emph{túshū} again
\end{itemize}

\begin{tcolorbox}[breakable,colback=teal!4,colframe=teal!35!black,title={Before you move on}]
What does \zh{图书} mean? (\textquotedblleft{}Books on the shelves (of a library or bookshop)\textquotedblright{}.) Say it once more.
\end{tcolorbox}
\section[\texorpdfstring{\zh{图} (tú)}{图 (tú)}]{\zh{图} — a frame, and five strokes inside it}

\label{lesson:ZH-W454-tu-chart}



\begin{quote}
The word first, then the character inside it.

\begin{itemize}
\raggedright
  \item \emph{Recall:} say \emph{túshū} — \textbf{R1}, one lesson back, before it decays
\end{itemize}

\begin{quote}
\zh{国} \zh{回}
\end{quote}

\emph{Read it:} each one aloud, with nothing said first — \emph{guó}, \emph{huí}
\end{quote}

\subsection*{Script — a frame, and five strokes inside it}
\hlblockfigure{\detokenize{figures/ZH-W454-tu-chart-filmstrip.pdf}}{How \zh{图} is written, stroke by stroke}

\begin{quote}
\zh{图}
\end{quote}

\textbf{A frame and what is inside it}, as in \zh{国}: the frame opens first, takes its inside, and closes last.

\textbf{The frame, three strokes.} A vertical down the left. A bar that turns down the right side. The bar across the foot comes last, after the inside.

\textbf{The inside, five strokes.} A short stroke falling to the left. A bar that turns down to the left. A long stroke falling to the right, across it. Then two dots, one under the other.

2 + 5 + 1 = eight.

\subsection*{Your turn}
\begin{itemize}
\raggedright
  \item \emph{Write it:} \zh{图}, all eight, with the model in front of you
  \item \emph{Cover:} write it again from memory, and count the pieces, not the strokes
\end{itemize}

\begin{tcolorbox}[breakable,colback=teal!4,colframe=teal!35!black,title={Before you move on}]
How many strokes in \zh{图}? (\textbf{Eight}.) Which stroke comes last? (\textbf{The bar that closes the frame at the foot}.) How many dots inside? (\textbf{Two}.)

Source: \href{https://raw.githubusercontent.com/chanind/hanzi-writer-data/68d10a4b21150cae5e1ebbd223eed289cf32d90c/data/\%E5\%9B\%BE.json}{Hanzi Writer Data \zh{图}, pinned stroke paths and medians}.
\end{tcolorbox}
\section[zhǐlù]{zhǐlù — to give directions; to show someone the way}

\label{lesson:ZH-C454-zhilu}



\begin{quote}
Before the new one: say the Mandarin for to make an inquiry, then the Mandarin for books on the shelves.
\end{quote}

\subsection*{You'll want to know: zhǐlù}
\emph{zhǐlù} — \textquotedblleft{}to give directions; to show someone the way\textquotedblright{}.

\begin{quote}
\emph{tā gěi wǒ zhǐlù gàosu wǒ wǎng qián zǒu} — He gave me directions and told me to go straight on.
\end{quote}

Its characters come later; for now it is a word for the ear.

\begin{cousinweb}
\emph{zhǐlù} is built from \emph{zhǐ} 'to point' and \zh{路} 'a road', both inherited from Middle Chinese, \emph{tsyijX} and \emph{luH} (Guangyun, in Baxter's transcription), which glosses the first as a finger and 'to show'.
\end{cousinweb}

\begin{grammarlens}[title={What you've built}]
One more word for getting around.
\end{grammarlens}

\subsection*{Your turn}
\begin{itemize}
\raggedright
  \item \emph{Say it:} \emph{zhǐlù}
  \item \emph{Say it:} \emph{zhǐlù}, once more
  \item \emph{Say it:} \emph{tā gěi wǒ zhǐlù gàosu wǒ wǎng qián zǒu} — He gave me directions and told me to go straight on.
  \item \emph{Recall:} say the Mandarin for to make an inquiry, then the Mandarin for books on the shelves, then say \emph{zhǐlù} again
\end{itemize}

\begin{tcolorbox}[breakable,colback=teal!4,colframe=teal!35!black,title={Before you move on}]
What does zhǐlù mean? (\textquotedblleft{}To give directions; to show someone the way\textquotedblright{}.) Say it once more.
\end{tcolorbox}
\section[xuānchuáncè]{xuānchuáncè — brochure; leaflet}

\label{lesson:ZH-C454-xuanchuance}



\begin{quote}
Before the new one: say the Mandarin for books on the shelves, then the Mandarin for to show someone the way.
\end{quote}

\subsection*{You'll want to know: xuānchuáncè}
\emph{xuānchuáncè} — \textquotedblleft{}brochure; leaflet\textquotedblright{}.

\begin{quote}
\emph{fúwùtái yǒu hěn duō miǎnfèi de xuānchuáncè} — The information desk has lots of free brochures.
\end{quote}

Its characters come later; for now it is a word for the ear.

\begin{cousinweb}
\emph{xuānchuáncè} is built from \emph{xuānchuán} 'to spread the word' and \emph{cè} 'a booklet', all inherited from Middle Chinese, \emph{sjwen}, \emph{drjwen} and \emph{tsrheak} (Guangyun, in Baxter's transcription), which glosses the first 'to spread', the second 'to turn' and the last as a set of bamboo slips.
\end{cousinweb}

\begin{grammarlens}[title={What you've built}]
One more word for getting around.
\end{grammarlens}

\subsection*{Your turn}
\begin{itemize}
\raggedright
  \item \emph{Say it:} \emph{xuānchuáncè}
  \item \emph{Say it:} \emph{xuānchuáncè}, once more
  \item \emph{Say it:} \emph{fúwùtái yǒu hěn duō miǎnfèi de xuānchuáncè} — The information desk has lots of free brochures.
  \item \emph{Recall:} say the Mandarin for books on the shelves, then the Mandarin for to show someone the way, then say \emph{xuānchuáncè} again
\end{itemize}

\begin{tcolorbox}[breakable,colback=teal!4,colframe=teal!35!black,title={Before you move on}]
What does xuānchuáncè mean? (\textquotedblleft{}Brochure; leaflet\textquotedblright{}.) Say it once more.
\end{tcolorbox}
